% Generated draft: requires mathematical and source review.
\documentclass[11pt]{article}
\usepackage[margin=25mm]{geometry}
\usepackage{fontspec}
\setmainfont{Latin Modern Roman}
\usepackage{amsmath,amssymb,mathtools,unicode-math}
\setmathfont{Latin Modern Math}
\usepackage{enumitem,needspace,xurl,hyperref}
\setlength{\parindent}{0pt}
\setlength{\parskip}{5pt}
\setlength{\emergencystretch}{3em}
\newcommand{\N}{\mathbb{N}}
\newcommand{\Z}{\mathbb{Z}}
\newcommand{\Q}{\mathbb{Q}}
\newcommand{\R}{\mathbb{R}}
\newcommand{\C}{\mathbb{C}}
\newcommand{\F}{\mathbb{F}}
\newcommand{\A}{\mathbb{A}}
\newcommand{\PP}{\mathbb{P}}
\newcommand{\E}{\mathbb{E}}
\newcommand{\eps}{\varepsilon}
\newcommand{\dd}{\,\mathrm{d}}
\begin{document}
\section*{UnsolvedMath ID: 20002670}
\textbf{Source problem code:} catalogue entry
\subsection{Definitions and mathematical statement}
Let \(H_\R\) have orthonormal basis \(e_1,\ldots,e_n\), and \(-1<q<1\). The \(q\)-Fock space is the vacuum space \(\C\Omega\) plus tensor powers, completed using the permutation-weighted inner product with weights \(q^{\operatorname{inv}(\sigma)}\). Write \(l_i\) for creation by \(e_i\), \(s_i=l_i+l_i^*\), and
\[
 \Gamma_{q,n}=W^*(s_1,\ldots,s_n),\qquad
 l_i^*l_j=q l_jl_i^*+\delta_{ij}I.
\]
The vacuum state is the usual tracial state. Question 10 contains four related tasks: determine dependence of the von Neumann algebra isomorphism type on \(q\), with combinatorial word recovery as a proposed approach; determine when generalized creation relations give a Cuntz extension; determine when generalized field algebras are type \(\mathrm{II}_1\) factors; and bound the completely bounded norms of length cutoffs. These tasks are not one solved assertion.

For word recovery, decompose a tensor word into maximal constant runs,
\[
 w=e_{i_1}^{\otimes\alpha_1}\otimes\cdots\otimes e_{i_k}^{\otimes\alpha_k},
 \quad\alpha_j\ge1,\quad i_j\ne i_{j+1},\quad N=\sum_j\alpha_j.
\]
Only adjacent run colors must differ; colors may repeat later. Define monic \(q\)-Hermite polynomials by
\[
 H_0^{(q)}=1,\quad H_1^{(q)}(x)=x,\quad
 H_{a+1}^{(q)}(x)=xH_a^{(q)}(x)-[a]_qH_{a-1}^{(q)}(x),
 \qquad[a]_q=\sum_{j=0}^{a-1}q^j.
\]
At \(q=0\) these are the monic scaled second-kind Chebyshev polynomials; naming them merely ``Chebytchev'' leaves a normalization ambiguity. Put \(v=H_{\alpha_1}^{(q)}(s_{i_1})\cdots H_{\alpha_k}^{(q)}(s_{i_k})\Omega\). If \(w=e_{j_1}\otimes\cdots\otimes e_{j_N}\) and \(\Pi_d\) is orthogonal projection onto tensor length \(d\), the aid's filtered formula is
\[
 \Pi_Nv=w,\qquad
 \Pi_{N-2}v=
 \sum_{\substack{1\le r<t\le N\\j_r=j_t\\r,t\text{ in distinct maximal runs}}}
 q^{t-r-1}\,
 e_{j_1}\otimes\cdots\widehat{e_{j_r}}\cdots\widehat{e_{j_t}}\cdots\otimes e_{j_N}.
\]
Hats denote deletion. Within-run contractions cancel. This identifies the first lower component, not all components of length \(N-4,N-6,\ldots\), nor an orthonormal tensor-word basis for \(q\ne0\).

The generalized relations in the source are
\[
 l_i^*l_j=\sum_{s,t}t_{i,j}^{s,t}l_sl_t^*+\delta_{ij}I.
\]
One must first specify coefficient conditions yielding a positive Fock representation and bounded operators. For such a representation the source asks when its creation \(C^*\)-algebra fits an exact sequence with compact ideal and Cuntz quotient, naturally \(0\to\mathcal K\to C^*(l_i)\to\mathcal O_n\to0\) in a finite \(n\)-generator interpretation, and when \(W^*(l_i+l_i^*)\) is a \(\mathrm{II}_1\) factor. No general admissibility or necessary-and-sufficient coefficient criterion is supplied. The aid reports factoriality for scalar \(q\) with finite \(n\ge2\); \(n=1\) is abelian and is excluded.

For the cutoff, let \(Q_d\) keep exact-length-\(d\) Wick words and kill other lengths, and \(P_{\le k}=\sum_{d=0}^kQ_d\). The source's \(P_k\) means length \emph{at most} \(k\), not exact length. The aid reports published bounds \(\|Q_d\|_{\mathrm{cb}}\le C_q(d+1)\), implying
\[
 \|P_{\le k}\|_{\mathrm{cb}}\le\frac{C_q(k+1)(k+2)}2.
\]
The constant may depend on \(q\). It reports the usual radial approximation consequence of weak-star complete metric approximation; the source additionally names CBAP of the noncommutative \(L^p\) spaces, \(1<p<\infty\), and exactness of the field \(C^*\)-algebra as intended consequences. Their precise approximation arguments and references are absent.

Finally the aid reports small-\(|q|\) isomorphism results for finite generators, unresolved dependence outside that regime, and known dependence in an infinite-generator analogue. Neither the smallness threshold nor the precise infinite-generator statement or citations are supplied. These status reports remain qualified; they do not follow from the word formula or cutoff estimate.
\subsection{Short English statement}
The record combines isomorphism, generalized creation/factoriality, word recovery, and length-cutoff questions. The aid reports a quadratic cutoff bound from exact-length estimates and gives a maximal-run \(q\)-Hermite formula for the leading and first lower tensor components. Other general questions remain separately qualified.
\subsection{Sources}
\begin{itemize}
\item Local registry, permanent UnsolvedMath ID \(20002670\), exact public code \texttt{catalogue entry}; canonical location \texttt{aim-probability-notes.json}, \texttt{notes[111]}.
\item AIM, \emph{Free Analysis: Problems}, Question 10, dated 24 August 2006, printed page 3, zero-based PDF page 2 according to the supplied indexed-text verification; reference \url{https://aimath.org/WWN/freeanalysis/freeanalysis.pdf}. No PDF was accessed here. No proposer, original remarks, or literature entry is supplied.
\item The aid reports published exact-length cb estimates, scalar factoriality, small-parameter isomorphisms, and an infinite-generator dependence result, but gives no authors, titles, theorem numbers, or complete bibliographic references.
\item Supplied research attempt 1: \texttt{partial\_result}, run and registry status \texttt{partially\_solved}, rigor label \texttt{proved}, low novelty confidence; classification \texttt{PARTIAL-PROGRESS}, created and updated 14 August 2026. The aid is machine generated, not peer reviewed, and requires independent verification before citation.
\end{itemize}
\end{document}
